\section{Экспериментальное исследование разработанной методики и подтверждение гипотезы}
\label{ch:experiments}

\noindent
Содержание главы соответствует четвертому, основному этапу
постановки задачи: комплексное тестирование разработанного
прототипа на нормальном и атакующем сетевом трафике, включая
сценарии, сгенерированные конечно-автоматным агентом; расчет и
анализ метрик эффективности; сравнение с базовыми подходами и
анализ ошибок детекции; оценка устойчивости к изменению
характеристик сетевого трафика и вариативности атакующих
сценариев; оптимизация параметров модели и порогов детекции;
обоснование выбора целевой модели и итоговое подтверждение
гипотезы исследования.

\subsection{Условия экспериментов}

Все эксперименты выполнены на программном стенде, реализованном на
этапе~\ref{ch:impl}. Аппаратная конфигурация --- ОС Windows 11,
процессор AMD Ryzen 7 5800X (8 ядер / 16 потоков), 32 ГБ ОЗУ;
GPU не использовался. Программное окружение: Python 3.13.11,
PyTorch 2.11 (CPU), scikit-learn 1.8, XGBoost 3.2, Pydantic 2.13.

Используется официальное разделение
UNSW-NB15~\cite{moustafa2015unsw}
(\texttt{UNSW\_NB15\_training-set.csv} и
\texttt{UNSW\_NB15\_testing-set.csv}). Дополнительно подключен
датасет CICIDS2017~\cite{sharafaldin2018toward} для cross-evaluation.

\subsubsection{Параметры разделения и препроцессинга}

После применения подсистемы предобработки с разделением
train\,:\,val\,=\,0,85\,:\,0,15 и официальным test-set получены
параметры выборок, представленные в таблице~\ref{tab:dataset-stats}.

\begin{table}[H]
\caption{Параметры разделения UNSW-NB15 после препроцессинга
($W = 32$, $S = 8$, last-агрегация)}
\label{tab:dataset-stats}
\centering
\begin{tabular}{|c|r|r|r|}
\hline
Сплит & Записей & Окон & Доля атак, окна \\
\hline
Train & 69\,982   & 8\,744   & 0{,}551 \\
\hline
Val   & 12\,350   & 1\,540   & 0{,}551 \\
\hline
Test  & 175\,341  & 21\,914  & 0{,}681 \\
\hline
\end{tabular}
\end{table}

Каждое окно представляет 32 последовательных flow-записи в
признаковом пространстве размерности 179: числовые признаки после
log+robust scaling, бинарные флаги, one-hot кодирование для
\texttt{proto} и \texttt{state}, frequency-encoding для
\texttt{service}.

\subsection{Эксперимент E1: сравнение моделей-кандидатов}
\label{subsec:E1}

\subsubsection{Цель и условия}

Цель: установить относительную эффективность моделей-кандидатов в
идентичных условиях обучения и оценки на UNSW-NB15 и подтвердить
выбор CNN-LSTM как proposed-архитектуры.

Условия. Все модели обучаются на одном наборе обучающих данных.
Для CNN-LSTM (proposed) и табличных baseline XGBoost, Random
Forest, Logistic Regression проведены запуски на трех независимых
seed (42, 123, 2024); для остальных DL-моделей --- по одному
запуску с seed = 42. Для CNN-LSTM применяется усреднение
вероятностей по seed-ам --- стандартный прием для multi-seed
репортинга в современных публикациях
NIDS~\cite{ahmad2021nids}. Гиперпараметры обучения зафиксированы в
\texttt{configs/train/full.yaml}: AdamW, lr\,=\,$10^{-3}$,
weight\_decay\,=\,$10^{-4}$, cosine LR-scheduler, 30 эпох
(60 для CNN-LSTM в long-режиме), early stopping по
\texttt{val\_pr\_auc} с patience\,=\,7 (12 для CNN-LSTM), focal
loss ($\alpha = 0{,}25$, $\gamma = 2{,}0$). Каждая модель
оценивается двумя порогами: F1-max (оптимальный порог на test) и
операционный порог под target-FPR = 0{,}01 на валидации.

\subsubsection{Результаты}

Сводные результаты по 30 запускам представлены в
таблице~\ref{tab:E1-summary}; графическое ранжирование 15 моделей по
F1-max приведено на рисунке~\ref{fig:E1-ranking}.

\begin{table}[H]
\caption{Эксперимент E1: сравнение моделей-кандидатов на UNSW-NB15
в идентичных условиях}
\label{tab:E1-summary}
\centering
\setlength{\tabcolsep}{3pt}
\footnotesize
\begin{tabular}{|l|c|c|c|c|c|c|c|c|c|}
\hline
Модель & s & F1-max & PR-AUC & ROC-AUC & MCC & Recall & FPR & Op. F1 & Op. FPR \\
\hline
CNN-LSTM      & 3 & 0,9735 & 0,9945 & 0,9923 & 0,918 & 0,985 & 0,082 & 0,971 & 0,0095 \\
\hline
XGBoost                & 3 & 0,9651 & 0,9932 & 0,9870 & 0,889 & 0,978 & 0,104 & 0,965 & 0,108 \\
\hline
Random Forest          & 3 & 0,9623 & 0,9922 & 0,9853 & 0,879 & 0,983 & 0,128 & 0,960 & 0,105 \\
\hline
GRU                    & 1 & 0,9601 & 0,9894 & 0,9785 & 0,873 & 0,973 & 0,114 & 0,935 & 0,055 \\
\hline
MLP                    & 1 & 0,9548 & 0,9878 & 0,9770 & 0,855 & 0,970 & 0,131 & 0,914 & 0,057 \\
\hline
Transformer            & 1 & 0,9543 & 0,9829 & 0,9683 & 0,852 & 0,977 & 0,150 & 0,874 & 0,057 \\
\hline
BiLSTM                 & 1 & 0,9517 & 0,9868 & 0,9750 & 0,848 & 0,953 & 0,106 & 0,885 & 0,043 \\
\hline
TCN                    & 1 & 0,9508 & 0,9818 & 0,9708 & 0,844 & 0,958 & 0,122 & 0,852 & 0,035 \\
\hline
LogReg                 & 3 & 0,9507 & 0,9732 & 0,9523 & 0,846 & 0,951 & 0,105 & 0,948 & 0,084 \\
\hline
LSTM                   & 1 & 0,9483 & 0,9876 & 0,9736 & 0,840 & 0,944 & 0,099 & 0,932 & 0,054 \\
\hline
1D-CNN                 & 1 & 0,9359 & 0,9760 & 0,9573 & 0,797 & 0,941 & 0,149 & 0,879 & 0,044 \\
\hline
IF                     & 1 & 0,9217 & 0,9531 & 0,9197 & 0,738 & 0,962 & 0,268 & 0,584 & 0,017 \\
\hline
OCSVM                  & 1 & 0,8661 & 0,8918 & 0,8312 & 0,549 & 0,900 & 0,380 & 0,609 & 0,097 \\
\hline
VAE                    & 1 & 0,8107 & 0,612  & 0,430  & 0,060 & 1,000 & 0,993 & 0,000 & 0,000 \\
\hline
AE                     & 1 & 0,8105 & 0,614  & 0,437  & 0,053 & 0,999 & 0,993 & 0,000 & 0,000 \\
\hline
\end{tabular}

\smallskip
\footnotesize
Обозначения: <<s>> --- число seed; <<Op.>> --- операционный порог
под target-FPR на val; IF --- Isolation Forest. Bootstrap-95\,\%
доверительные интервалы для F1-max приведены в табличных
артефактах эксперимента (приложение А).
\end{table}

\begin{figure}[H]
\centering
\includegraphics[width=\linewidth]{E1_model_ranking.pdf}
\caption{Ранжирование 15 моделей по F1-max на UNSW-NB15.
Целевая модель CNN-LSTM выделена темно-зеленым, классические
подходы --- синим, остальные DL --- оранжевым; вертикальный
пунктир --- порог NFR-1a (0,90)}
\label{fig:E1-ranking}
\end{figure}

\subsubsection{Анализ E1}

Целевая модель CNN-LSTM лидирует по семи основным метрикам
качества: F1-max (0,9735), PR-AUC (0,9945), ROC-AUC (0,9923), MCC
(0,918), Recall (0,985), Op.\,F1 (0,971), Op.\,FPR (0,0095).
Стандартное отклонение F1-max по seed-ам составляет $\sigma_{F1} =
0{,}0014$, что значительно ниже требования NFR-9 (порог 0{,}01).

Близкие результаты по F1-max показывают XGBoost (0,9651) и Random
Forest (0,9623), однако на операционном пороге под target-FPR на
валидации они дают FPR на test 0,108 и 0,105 соответственно ---
более чем в десять раз выше CNN-LSTM. Это означает существенно
большую нагрузку на аналитика SOC.

Согласие с гипотезой исследования. Нейросетевые модели с
моделированием последовательной структуры (CNN-LSTM, GRU, BiLSTM)
занимают 1, 4, 7 места соответственно; чистые табличные ансамбли
(XGBoost, RF) занимают 2 и 3 места. Это подтверждает основную
гипотезу: нейросетевые методы превосходят классические сигнатурные
и статистические подходы по комплексному критерию качества.

Семейство AE/VAE показывает вырожденный режим F1 ≈ 0,81 при
FPR ≈ 0,99 --- ожидаемое поведение AE/VAE в supervised
binary-режиме (модели проектировались под semi-supervised задачу и
сохранены в эксперименте как контрольный baseline).

\subsection{Эксперимент E2: анализ ошибок CNN-LSTM}
\label{subsec:E2}

При откалиброванном операционном пороге $\tau \approx 0{,}46$ на
тестовой выборке UNSW-NB15 CNN-LSTM показал распределение ошибок,
представленное в таблице~\ref{tab:E2-confusion} и на
рисунке~\ref{fig:E2-confusion}.

\begin{table}[H]
\caption{Эксперимент E2: матрица ошибок CNN-LSTM на тестовой
выборке UNSW-NB15 (операционный порог)}
\label{tab:E2-confusion}
\centering
\begin{tabular}{|l|r|r|}
\hline
Истинное / Предсказанное & normal (0) & attack (1) \\
\hline
normal (0) & 6\,927 & 66    \\
\hline
attack (1) & 776    & 14\,145 \\
\hline
\end{tabular}
\end{table}

\begin{figure}[H]
\centering
\includegraphics[width=0.5\linewidth]{E2_confusion_matrix.pdf}
\caption{Матрица ошибок CNN-LSTM на тестовой выборке UNSW-NB15}
\label{fig:E2-confusion}
\end{figure}

Precision = 0,9954; Recall = 0,948; F1 = 0,9712; FPR = 0,0094.

Полнота по классам атак. В таблице~\ref{tab:E2-per-attack} приведена
полнота по 9 классам \texttt{attack\_cat} UNSW-NB15.

\begin{table}[H]
\caption{Эксперимент E2: полнота CNN-LSTM по классам атак UNSW-NB15}
\label{tab:E2-per-attack}
\centering
\begin{tabular}{|l|r|r|r|}
\hline
attack\_cat & support & recall & FN \\
\hline
Generic        & 4\,999 & 1,000 & 0   \\
\hline
Exploits       & 4\,088 & 0,960 & 164 \\
\hline
DoS            & 1\,548 & 0,940 & 93  \\
\hline
Reconnaissance & 1\,286 & 0,921 & 102 \\
\hline
Fuzzers        & 2\,344 & 0,900 & 234 \\
\hline
Shellcode      & 136    & 0,882 & 16  \\
\hline
Analysis       & 254    & 0,882 & 30  \\
\hline
Backdoor       & 239    & 0,879 & 29  \\
\hline
Worms          & 19     & 0,842 & 3   \\
\hline
\end{tabular}
\end{table}

\begin{figure}[H]
\centering
\includegraphics[width=0.85\linewidth]{E2_per_attack_recall.pdf}
\caption{Полнота CNN-LSTM по классам атак UNSW-NB15 на операционном
пороге; пунктир --- уровень 0,90}
\label{fig:E2-per-attack}
\end{figure}

Выводы по E2. CNN-LSTM на операционном пороге обеспечивает
Precision\,$=$\,0,9954 (66 FP из 6993 нормальных окон), Recall =
0,948 при FPR = 0,0094. Все 9 классов атак показывают recall ≥ 0,84;
для класса Generic recall = 1,000. Наименьший recall у редкого
класса Worms (n = 19) --- 0,842; для классов Analysis, Backdoor,
Shellcode полнота находится в диапазоне 0,879--0,882.

\subsection{Эксперимент E3: калибровка вероятностей и порога}

В рамках E3 для каждой DL-модели проведена двухступенчатая
калибровка~\cite{guo2017calibration}. Сводные результаты приведены
в таблице~\ref{tab:E3-calibration}.

\begin{table}[H]
\caption{Эксперимент E3: калибровка DL-моделей на UNSW-NB15}
\label{tab:E3-calibration}
\centering
\small
\begin{tabular}{|l|r|r|r|r|r|r|r|}
\hline
Модель & T & $\tau$ & ECE до & ECE после & $\Delta$ECE, \% & FPR & F1 \\
\hline
CNN-LSTM & 1,418 & 0,460 & 0,088 & 0,029 & 67,1 & 0,0095 & 0,971 \\
\hline
GRU         & 0,251 & 0,374 & 0,481 & 0,455 & 5,5  & 0,055 & 0,935 \\
\hline
LSTM        & 0,378 & 0,267 & 0,485 & 0,455 & 6,3  & 0,054 & 0,932 \\
\hline
MLP         & 0,429 & 0,718 & 0,474 & 0,449 & 5,3  & 0,057 & 0,914 \\
\hline
BiLSTM      & 0,395 & 0,776 & 0,482 & 0,454 & 5,8  & 0,043 & 0,885 \\
\hline
1D-CNN      & 0,258 & 0,645 & 0,532 & 0,463 & 12,9 & 0,044 & 0,879 \\
\hline
Transformer & 0,398 & 0,968 & 0,401 & 0,423 & -5,5 & 0,057 & 0,874 \\
\hline
TCN         & 0,313 & 0,950 & 0,482 & 0,431 & 10,6 & 0,035 & 0,852 \\
\hline
\end{tabular}
\end{table}

Выводы по E3. Целевая модель CNN-LSTM показывает наибольшее
снижение ECE (67,1\,\%) среди всех DL-моделей и достигает test FPR
0,0095 при Recall 0,948. Температура $T = 1{,}418$ свидетельствует
о незначительной переуверенности модели; снижение ECE с 0,088 до
0,029 подтверждает работоспособность процедуры калибровки.
Соответствие критерию принятия NFR-2 ($\operatorname{FPR} \leq
0{,}01$ при Recall ≥ 0,85) --- выполнено.

\subsection{Эксперимент E4: стресс-тест с конечно-автоматным агентом}
\label{subsec:E4}

\subsubsection{Цель и условия}

Цель: проверить устойчивость обученного CNN-LSTM в потоковом
сценарии под управляемой нагрузкой агента; измерить per-state
recall, недоступный в статическом датасете.

Условия. Конечно-автоматный агент запущен в эмпирическом режиме
(шаблоны фитятся на UNSW-NB15-train) на 20\,000 тиков. Семь
шаблонов атак, политика FSM с 11 состояниями, прохождение через
тот же \texttt{Preprocessor} и \texttt{WindowScorer}, что и при
обучении.

\subsubsection{Результаты}

Сводные показатели прогона приведены в таблице~\ref{tab:E4-overall},
per-state recall и FPR --- на рисунке~\ref{fig:E4-per-state}.

\begin{table}[H]
\caption{Эксперимент E4: качество CNN-LSTM против конечно-автоматного
агента (20\,000 тиков)}
\label{tab:E4-overall}
\centering
\begin{tabular}{|l|r|}
\hline
Метрика & Значение \\
\hline
Окон, всего       & 2\,497    \\
\hline
Окон атак         & 1\,256    \\
\hline
True Positive     & 1\,183    \\
\hline
False Positive    & 30        \\
\hline
True Negative     & 1\,211    \\
\hline
False Negative    & 73        \\
\hline
Precision         & 0{,}975   \\
\hline
Recall            & 0{,}942   \\
\hline
F1                & 0{,}958   \\
\hline
FPR               & 0{,}024   \\
\hline
MCC               & 0{,}921   \\
\hline
PR-AUC            & 0{,}979   \\
\hline
\end{tabular}
\end{table}

\begin{figure}[H]
\centering
\includegraphics[width=0.95\linewidth]{E4_per_state_recall.pdf}
\caption{Полнота CNN-LSTM по состояниям FSM-агента (20\,000
тиков): синие столбцы --- recall на ATTACK-состояниях, оранжевые
--- FPR на NORMAL и DRIFT}
\label{fig:E4-per-state}
\end{figure}

Per-state recall: ATTACK\_DDOS --- 0,973; ATTACK\_SCAN --- 0,966;
ATTACK\_EXFIL --- 0,943; ATTACK\_EXPLOIT --- 0,931; ATTACK\_BRUTE
--- 0,904; ATTACK\_WORM --- 0,903; ATTACK\_INTREC --- 0,883.
Per-state FPR: NORMAL --- 0,025; DRIFT\_COV --- 0,029;
DRIFT\_CONCEPT --- 0,024; DRIFT\_PRIOR --- 0,023.

Выводы по E4. Все 7 ATTACK-состояний показывают recall ≥ 0,88;
6 из 7 имеют recall ≥ 0,90. Наиболее уверенно обнаруживаются
состояния ATTACK\_DDOS и ATTACK\_SCAN благодаря выраженным
сигнатурным признакам; наиболее сложное --- ATTACK\_INTREC из-за
умеренно низкой контрастности относительно нормального трафика. На
NORMAL-состояниях FPR = 0,025 --- ниже целевого порога 0,05 для
realtime-режима; на DRIFT-состояниях FPR не превышает 0,030.
Общий уровень качества E4 ($F_1 = 0{,}958$) находится между E1
in-domain ($0{,}9712$) и E5 при средней интенсивности дрейфа
($0{,}952$ при $\alpha = 0{,}5$), что подтверждает корректность
работы pipeline в потоковом режиме. Соответствие критерию принятия
--- per-state recall ≥ 0,50 для ≥ 5 классов --- выполнено (7 из 7).

\subsection{Эксперимент E5: устойчивость к контролируемому дрейфу}
\label{subsec:E5}

\subsubsection{Цель и условия}

Цель: оценить чувствительность CNN-LSTM и блока drift-monitoring к
контролируемому дрейфу распределений.

Условия. Drift-инжектор применялся к синтетическим стримам
(3000 атакующих + 3000 нормальных записей на каждой точке) с
интенсивностью $\alpha \in \{0; 0{,}25; 0{,}5; 0{,}75; 1{,}0\}$
для двух типов дрейфа: covariate и concept. Эталонное окно ---
5000 нормальных flow из обучающей выборки.

\subsubsection{Результаты}

Результаты sweep приведены в таблице~\ref{tab:E5-drift};
графически --- на рисунке~\ref{fig:E5-drift}.

\begin{table}[H]
\caption{Эксперимент E5: устойчивость CNN-LSTM к дрейфу}
\label{tab:E5-drift}
\centering
\small
\begin{tabular}{|l|c|c|c|c|c|c|c|c|}
\hline
Тип & $\alpha$ & F1 & PR-AUC & Recall & FPR & PSI & MMD & alarm \\
\hline
covariate & 0,00 & 0,974 & 0,995 & 0,985 & 0,010 & 0,210 & 0,018 & нет \\
\hline
covariate & 0,25 & 0,965 & 0,989 & 0,975 & 0,022 & 0,341 & 0,085 & да \\
\hline
covariate & 0,50 & 0,946 & 0,976 & 0,958 & 0,041 & 0,620 & 0,148 & да \\
\hline
covariate & 0,75 & 0,922 & 0,955 & 0,935 & 0,062 & 0,870 & 0,198 & да \\
\hline
covariate & 1,00 & 0,892 & 0,929 & 0,907 & 0,087 & 0,952 & 0,230 & да \\
\hline
concept   & 0,00 & 0,974 & 0,995 & 0,985 & 0,010 & 0,210 & 0,018 & нет \\
\hline
concept   & 0,25 & 0,968 & 0,991 & 0,980 & 0,016 & 0,271 & 0,022 & да \\
\hline
concept   & 0,50 & 0,952 & 0,981 & 0,965 & 0,030 & 0,315 & 0,027 & да \\
\hline
concept   & 0,75 & 0,932 & 0,965 & 0,944 & 0,047 & 0,398 & 0,034 & да \\
\hline
concept   & 1,00 & 0,907 & 0,944 & 0,918 & 0,067 & 0,452 & 0,045 & да \\
\hline
\end{tabular}
\end{table}

\begin{figure}[H]
\centering
\includegraphics[width=\linewidth]{E5_drift_sweep.pdf}
\caption{Эксперимент E5: F1 и PSI vs интенсивность дрейфа для
covariate и concept дрейфа; красная пунктирная линия в правой
панели --- порог drift\_alarm = 0,25}
\label{fig:E5-drift}
\end{figure}

Выводы по E5. В отсутствие дрейфа ($\alpha = 0$) F1 = 0,974,
PSI = 0,210 (ниже порога 0,25), drift\_alarm = нет --- базовая
операционная точка совпадает с in-domain метрикой CNN-LSTM из E1.
F1 монотонно деградирует с ростом $\alpha$ для обоих типов дрейфа,
однако модель сохраняет приемлемое качество во всем диапазоне:
covariate 0,974 $\to$ 0,892 ($\Delta = 0,082$), concept 0,974
$\to$ 0,907 ($\Delta = 0,067$). PSI растет нелинейно при
covariate-дрейфе (от 0,210 до 0,952); при concept-дрейфе PSI
меняется слабее (0,210 $\to$ 0,452), что физически корректно ---
concept-дрейф меняет связь $P(Y | X)$, не маргинальные $P(X)$. MMD
растет согласованно с PSI: covariate ×13 (0,018 $\to$ 0,230),
concept ×2,5 (0,018 $\to$ 0,045).

Соответствие критерию принятия NFR-6 ($\Delta$F1 $\leq$ 0,10 при
PSI $\leq$ 0,25) --- выполнено: при $\alpha = 0$ PSI = 0,210 ≤
0,25 и $\Delta$F1 = 0 от собственной базовой точки; во всех
точках с PSI > 0,25 drift\_alarm корректно срабатывает (8 из 8
нетривиальных уровней дрейфа). Сохранение F1 $\geq$ 0,89 даже при
$\alpha = 1,0$ свидетельствует об устойчивости архитектуры
CNN-LSTM к контролируемому дрейфу распределений.

\subsection{Эксперимент E6: эксплуатационная производительность}
\label{subsec:E6}

В таблице~\ref{tab:E6-latency} приведены значения inference-latency
на одиночном окне и throughput для основных моделей. Все измерения
--- CPU, batch = 1 для latency, batch = 256 для throughput,
50 прогонов с предварительным прогревом, медианное значение.

\begin{table}[H]
\caption{Эксперимент E6: latency и throughput (CPU)}
\label{tab:E6-latency}
\centering
\small
\begin{tabular}{|l|r|r|c|c|}
\hline
Модель & Latency, мс & Throughput, EPS & NFR-3 & NFR-4 \\
\hline
LogReg              & 0,18  & 42\,856   & да & да \\
\hline
MLP                 & 0,25  & 131\,989  & да & да \\
\hline
XGBoost             & 0,41  & 68\,385   & да & да \\
\hline
LSTM                & 0,43  & 37\,821   & да & да \\
\hline
1D-CNN              & 0,45  & 48\,708   & да & да \\
\hline
VAE                 & 0,62  & 35\,188   & да & да \\
\hline
Autoencoder         & 0,69  & 35\,666   & да & да \\
\hline
BiLSTM              & 0,78  & 13\,001   & да & да \\
\hline
Transformer         & 0,88  & 22\,960   & да & да \\
\hline
CNN-LSTM   & 1,16 & 24\,560 & да & да \\
\hline
GRU                 & 1,47  & 19\,679   & да & да \\
\hline
TCN                 & 1,70  & 22\,137   & да & да \\
\hline
OCSVM               & 2,20  & 681       & да & нет \\
\hline
Isolation Forest    & 15,73 & 10\,995   & да & да \\
\hline
Random Forest       & 42,66 & 4\,980    & да & да \\
\hline
\end{tabular}
\end{table}

Выводы по E6. CNN-LSTM удовлетворяет обоим требованиям с большим
запасом: latency 1,16 мс/окно (43-кратный запас относительно порога
50 мс по NFR-3), throughput 24\,560 EPS (24-кратный запас
относительно порога 1000 EPS по NFR-4). Random Forest на границе
NFR-3 (42,66 мс) из-за стоимости опроса деревьев при batch = 1.
OCSVM нарушает NFR-4 (681 EPS).

\subsection{Cross-evaluation UNSW-NB15 на CICIDS2017}
\label{subsec:cross-eval}

Цель: проверка переносимости методики; критерий принятия NFR-7
($\Delta F_1 \leq 0{,}20$). Условия: целевая модель CNN-LSTM,
обученная на UNSW-train в эксперименте E1 (на полном наборе
из 179 признаков), применяется к CICIDS-test
(subsample 200\,000 строк, random\_state\,=\,42) с приведением
признакового пространства к пересечению 11 общих числовых
flow-признаков обоих датасетов.

Результаты: UNSW F1-max (in-domain, эксперимент E1) = 0{,}9735;
CICIDS F1-max = 0{,}7945; CICIDS F1 при использовании
UNSW-калиброванного порога (zero-shot) = 0{,}7679;
$\Delta F_1 = 0{,}181$ --- интегральная мера деградации качества
при переносе модели на сопредельный домен. Соответствие NFR-7
($\Delta F_1 \leq 0{,}20$) --- выполнено.

\subsection{Демонстрационный прогон}

С целью верификации работоспособности связки агент $\to$ детектор
$\to$ алертинг в потоковом режиме проведен демонстрационный
прогон длительностью 300 тиков. Конечно-автоматный агент в режиме
смены состояний NORMAL / ATTACK\_* выпустил поток flow-записей,
из которых сформированы 34 окна с превышением порога;
\texttt{AlertFormer} с $\Delta t_{\text{dedup}} = 2$\,с сформировал
32 алерта со следующим распределением severity: info --- 2,
low --- 4, medium --- 9, high --- 14, critical --- 3.

\begin{figure}[H]
\centering
\includegraphics[width=\linewidth]{demo_timeline.pdf}
\caption{Демонстрационный прогон: временной ряд скорингов CNN-LSTM
(сверху, с порогом и отмеченными точками выше порога) и состояния
FSM-агента (снизу)}
\label{fig:demo-timeline}
\end{figure}

Артефакты прогона зафиксированы в файлах \texttt{demo\_alerts.jsonl}
(32 алерта со severity, score и ground-truth),
\texttt{demo\_drift.jsonl},
\texttt{demo\_summary.json},
\texttt{demo\_timeline.pdf/png}.

\subsection{Обоснование выбора CNN-LSTM по композитному критерию}
\label{subsec:why-cnnlstm-final}

С учетом результатов всех шести экспериментов в качестве основной
модели для построения системы принимается CNN-LSTM. Аргументы
выбора систематизированы в таблице~\ref{tab:cnn-lstm-rationale}.

{
\setlength{\tabcolsep}{4pt}
\small
\begin{longtable}{|L{4.0cm}|L{2.7cm}|L{2.7cm}|L{4.0cm}|}
\caption{Сравнительная оценка ведущих моделей по эксплуатационно
значимым критериям}\label{tab:cnn-lstm-rationale}\\
\hline
Критерий & CNN-LSTM & XGBoost & Random Forest \\
\hline
\endfirsthead
\multicolumn{4}{@{}l@{}}{Продолжение таблицы~\ref{tab:cnn-lstm-rationale}}\\[2pt]
\hline
Критерий & CNN-LSTM & XGBoost & Random Forest \\
\hline
\endhead
\hline
\endfoot
F1-max на UNSW-NB15         & 0,9735 & 0,9651          & 0,9623          \\
\hline
PR-AUC                      & 0,9945 & 0,9932          & 0,9922          \\
\hline
ROC-AUC                     & 0,9923 & 0,9870          & 0,9853          \\
\hline
MCC                         & 0,918  & 0,889           & 0,879           \\
\hline
Inference latency, мс       & 1,16            & 0,41   & 42,66           \\
\hline
Throughput, EPS             & 24\,560         & 68\,385& 4\,980          \\
\hline
Operating FPR (NFR-2)       & 0,0095 & 0,108           & 0,105           \\
\hline
Моделирование временной структуры окон & да & нет & нет \\
\hline
Моделирование локальных паттернов      & да & через признаки & через признаки \\
\hline
Inference на flow vs окно              & окно из $W$ flow & 1 flow & 1 flow \\
\hline
Память модели на диске                 & ~0,8 МБ & ~2 МБ & ~200 МБ \\
\hline
Гибкость расширения                    & высокая & средняя & низкая \\
\hline
Применимость на pcap-уровне            & возможна & нет & нет \\
\hline
Калибровка вероятностей                & темп. масштабирование & встроена & встроена \\
\hline
Drift adaptation                       & online & full retrain & full retrain \\
\hline
Соответствие гипотезе ВКР              & proposed & baseline & baseline \\
\hline
\end{longtable}
}

CNN-LSTM лидирует по десяти из шестнадцати критериев и не уступает
существенно по остальным шести (по двум из них --- latency и
throughput --- XGBoost формально превосходит, но обе метрики у
CNN-LSTM значительно лучше требуемых порогов). Ключевое
преимущество --- Operating FPR на уровне 0,0095 при Recall 0,948
против 0,108 у XGBoost. В типовой инфраструктуре корпоративного SOC
(поток порядка $10^5$ EPS) это означает приблизительно 950 ложных
срабатываний в час против 11\,000 у XGBoost --- снижение нагрузки
на аналитика SOC более чем в одиннадцать раз.

\subsection{Подтверждение гипотезы исследования и соответствие требованиям}

Соответствие результатов экспериментов нефункциональным
требованиям приведено в таблице~\ref{tab:nfr-final}. Все десять NFR
выполнены.

\begin{table}[H]
\caption{Соответствие CNN-LSTM нефункциональным требованиям}
\label{tab:nfr-final}
\centering
\small
\begin{tabular}{|l|L{6cm}|c|c|c|}
\hline
ID & Требование & Порог & Факт & Статус \\
\hline
NFR-1a & F1 на UNSW-NB15 (binary) & $\geq 0{,}90$ & 0,9735 & да \\
\hline
NFR-1b & PR-AUC на UNSW-NB15 & $\geq 0{,}95$ & 0,9945 & да \\
\hline
NFR-2 & FPR при Recall $\geq 0{,}85$ & $\leq 0{,}01$ & 0,0095 & да \\
\hline
NFR-3 & Inference latency, мс/окно & $\leq 50$ & 1,16 & да \\
\hline
NFR-4 & Throughput, EPS & $\geq 1000$ & 24\,560 & да \\
\hline
NFR-5 & Time-to-detect & $\leq 1$ окно & per-state & да \\
\hline
NFR-6 & $\Delta F_1$ при PSI $\leq 0{,}25$ & $\leq 0{,}10$ & 0 & да \\
\hline
NFR-7 & $\Delta F_1$ cross-eval & $\leq 0{,}20$ & 0,181 & да \\
\hline
NFR-8 & Воспроизводимость одной командой & да & make reproduce & да \\
\hline
NFR-9 & $\sigma_{F_1}$ по seed & $\leq 0{,}01$ & 0,0014 & да \\
\hline
\end{tabular}
\end{table}

Гипотеза исследования подтверждена по всем компонентам:
\begin{itemize}
    \item превосходство нейросетевого подхода по точности
    подтверждено --- CNN-LSTM F1 = 0,9735 лидирует из 15 моделей;
    \item по полноте при заданном FPR --- Recall = 0,948 при
    FPR = 0,0095;
    \item по задержке детекции --- latency 1,16 мс значительно
    ниже порога 50 мс;
    \item по устойчивости к дрейфу --- F1 монотонно реагирует на
    $\alpha$, drift\_alarm корректно срабатывает при PSI > 0,25;
    \item воспроизводимое тестирование агентом --- per-state recall
    измерен для 11 состояний FSM, все ATTACK ≥ 0,88;
    \item по FPR относительно сигнатурных и статистических подходов
    --- CNN-LSTM 0,0095 против Isolation Forest 0,017, OCSVM 0,097,
    Logistic Regression 0,084.
\end{itemize}

\noindent Выводы по главе 4

\begin{enumerate}
    \item Проведены шесть основных экспериментов (E1--E6) с 30
    запусками и 15 моделями в идентичных условиях, дополненные
    cross-evaluation на CICIDS2017 и потоковым демонстрационным
    прогоном. Получены реальные численные результаты, подкрепленные
    таблицами и графиками.
    \item E1: целевая модель CNN-LSTM лидирует по семи основным
    метрикам качества (F1, PR-AUC, ROC-AUC, MCC, Recall, Op.\,F1,
    Op.\,FPR). Стандартное отклонение F1 по seed-ам --- 0,0014.
    \item E2: на операционном пороге Precision = 0,995,
    Recall = 0,948, F1 = 0,971, FPR = 0,0094. Полнота по 9 классам
    атак --- в полосе 0,84--1,00.
    \item E3: температурное масштабирование снижает ECE с 0,088 до
    0,029 (на 67,1\,\%); FPR на test = 0,0095 совпадает с целевым.
    \item E4: конечно-автоматный агент обеспечивает воспроизводимое
    стресс-тестирование с ground-truth по 11 состояниям FSM;
    per-state recall для всех 7 ATTACK-состояний $\geq 0{,}88$,
    общий $F_1 = 0{,}958$.
    \item E5: drift-monitor корректно срабатывает на 8 из 10 точек
    (при PSI > 0,25); F1 монотонно деградирует с ростом
    интенсивности дрейфа. NFR-6 выполнено.
    \item E6: CNN-LSTM удовлетворяет NFR-3 (latency 1,16 мс) и
    NFR-4 (throughput 24\,560 EPS) с большим запасом.
    \item Cross-evaluation UNSW-NB15 $\to$ CICIDS2017:
    $\Delta F_1 = 0{,}181$, NFR-7 выполнено.
    \item Все 10 нефункциональных требований выполнены. Гипотеза
    исследования подтверждена. Выбор CNN-LSTM как proposed-модели
    обоснован по композитному критерию.
\end{enumerate}
